\section{One World out of Many: Collective Geometry, Observable Sectors, and Observer-Dependent Reality}\label{sec:collective_geometry}

Wheeler's third question asks how ``one world'' arises out of many observer-participants, and in answering it he adopts F{\o}llesdal's dictum that meaning is ``the joint product of all the evidence that is available to those who communicate''~\cite{Wheeler1990}; a claimed measurement, he adds, counts only once it can be checked with one's fellows. This section develops the paper's formal counterpart of that strand, item W3 of Table~\ref{tab:wheeler_correspondence}. The agents of Part~I communicate only through the transported relative entropies $\beta_{ij} D_{\mathrm{KL}}(q_i \| \Omega_{ij}[q_j])$ of the consensus functional, Eq.~\eqref{eq:free_energy_functional_final}, and each agent also carries its own observation term. At a stationary point the attention rows are the Gibbs rows of Section~\ref{sec:complete_free_energy}, and for fixed neighbors each agent's optimal belief is the normalized geometric mean of its prior and the transported beliefs it attends to (Theorem~\ref{thm:representation_app}), reweighted by a power of the agent's own likelihood; for Gaussian beliefs and a linear Gaussian likelihood it is a precision-weighted average. The residual transported divergence is the frame-independent part of the disagreement that remains. The stationarity conditions are derived (status D); reading the stationary point as ``one world'' is an interpretation (status S; Table~\ref{tab:claim_status}). The model has the structure Wheeler names, a shared description assembled from the evidence available to agents who communicate, and it departs from his proposal in three ways: the evidence is continuous, the channel is a declared divergence, and agreement is graded, never a registered yes-no fact.

\subsection{The Consensus Metric and Its Regulator Obstruction}
\label{sec:consensus_metric}

Each agent induces its own metric on the base manifold, and the question is how these individual geometries combine into shared structure. The answer depends on which block of the combined tensor of Eq.~\eqref{eq:combined_tensor} is averaged. The Fisher block $G^{(q)}_i$ of Eq.~\eqref{eq:induced_metric_full}, computed with the covariant jet of Eq.~\eqref{eq:covariant_jet}, is the ordinary pullback of the frame-stripped law $\hat q_i = \rho(U_i^{-1}) q_i$, and a passive local re-expression $U_i \mapsto g_i(c) U_i$, $q_i \mapsto \rho(g_i(c)) q_i$ leaves it unchanged (Section~\ref{sec:pullback} and~\cite[Ch.~B2]{DennisGaugeVFESuppB2026}). The ordinary Fisher block of $q_i$ is likewise unchanged by the change-of-reference action $U_i \mapsto U_i g(c)$ applied to all agents, which moves neither beliefs nor transports. The presence-weighted consensus metric
\begin{equation}
\bar{G}^{\mathrm{cons}}_{\mu\nu}(c) = \sum_i w_i(c) G^{(q)}_{i,\mu\nu}(c)
\label{eq:consensus_metric_covariant}
\end{equation}
is therefore gauge invariant with no orbit average, although it is relative to the agents' connections. It is not the only invariant. Relative-entropy values, the conjugacy classes of loop holonomies, and the per-edge twists modulo one simultaneous conjugation survive every change of frame~\cite[Ch.~14]{DennisGaugeVFE2026}, and the Fisher clock of a natural-gradient history of an invariant free energy equals its gauge-quotient duration~\cite[Ch.~7]{DennisGaugeVFE2026}. An orbit average is needed only for the frame-twist block $\mathrm{tr}(A^{(i)}_\mu A^{(i)}_\nu)$, with $A^{(i)} = U_i^{-1} dU_i$, which is where the signature constructions of Section~\ref{sec:signature_resolution} live, and there it meets the following obstruction.
\begin{proposition}[Regulator obstruction for the frame-twist block]\label{prop:regulator_obstruction}
Let $T_{i,\mu\nu} = \mathrm{tr}(A^{(i)}_\mu A^{(i)}_\nu)$.
(a) $T_i$ is invariant under constant right translation $U_i \mapsto U_i g$, so a Haar average over one copy of $G$ is redundant.
(b) Under the c-dependent right translation $U_i \mapsto U_i g(c)$ of all agents, which leaves every transport $\Omega_{ij}$ and every belief unchanged, and hence every alignment, attention, and cross-scale term, $A^{(i)} \mapsto g^{-1} A^{(i)} g + g^{-1} dg$, and the choice $g = U_i^{-1}$ gives $A^{(i)} = 0$ and $T_i = 0$ for any one chosen agent. The frame-twist block has no gauge-invariant content.
(c) After discretization to finitely many base points, the orbit average over c-dependent $g$ exists for compact $G$, since the product Haar measure on $G^N$ is a probability measure. It has no continuum limit: the group of maps $\mathcal{C} \to G$ is infinite-dimensional and not locally compact, so the Haar construction does not apply, and at lattice spacing $\varepsilon$ the discrete Maurer-Cartan term between independent Haar-distributed neighbors is of order $\varepsilon^{-1}$, so the average of $T_i$ grows in magnitude as $\varepsilon^{-2}$.
(d) For noncompact $G$, including $\mathrm{GL}(K_q)$ and $\mathrm{SO}(1,3)$, the Haar measure has infinite total mass already for a single copy, and the average is undefined without a further regulator.
\end{proposition}
\begin{proof}
(a) $\mathrm{tr}(g^{-1}A_\mu g g^{-1}A_\nu g) = \mathrm{tr}(A_\mu A_\nu)$ by cyclicity; this is the residual-subgroup statement of Section~\ref{sec:dual_role_rigor}. (b) $(U_i g)^{-1} d(U_i g) = g^{-1}(U_i^{-1} dU_i) g + g^{-1} dg$ and $(U_i g)(U_j g)^{-1} = U_i U_j^{-1}$; with $g = U_i^{-1}$ the transformed frame of agent $i$ is the identity at every $c$. (c) A product of finitely many compact groups is compact, and its normalized Haar measure is the product measure. For independent Haar-distributed $g_x, g_{x+\varepsilon}$ the element $g_x^{-1}g_{x+\varepsilon}$ is Haar distributed and its logarithm is of order one almost surely, so the lattice connection $\varepsilon^{-1}\log(g_x^{-1}g_{x+\varepsilon})$ is of order $\varepsilon^{-1}$ and its trace square, which is negative definite on the compact algebra, has expectation proportional to $\varepsilon^{-2}$. (d) A noncompact locally compact group has infinite Haar measure.
\end{proof}
In the flat Regime~I every curvature invariant vanishes, since $F = 0$. With independent links on the interaction graph (Regime~II) the invariants are the conjugacy classes of loop holonomies listed above, and their operational detection is item C1 of Section~\ref{sec:open_problems}. The construction transfers verbatim from the state fiber to the model fiber.

\subsection{Observable, Subthreshold, and Internal Sectors}
\label{sec:observable_sectors}

The covariant Fisher block $G_i(c) := G^{(q)}_i(c)$ is a positive semidefinite tensor on $T_c\mathcal{C}$ whose rank equals the rank of $d\sigma_i(c)$ and is therefore at most $\dim\mathcal{C}$. Its eigendirections split $T_c\mathcal{C}$ into three sectors against thresholds: an observable sector $\mathcal{D}_{\text{obs}}$ of directions with large eigenvalue, read as perceived spacetime, a subthreshold sector $\mathcal{D}_{\text{subthresh}}$ carrying non-negligible but imperceptible information rate, and an internal sector $\mathcal{D}_{\text{internal}}$ along which the section is nearly constant, with $\lambda_{\text{obs}} \gg \lambda_{\text{subthresh}} \gg \lambda_{\text{internal}} \approx 0$ assumed to hold smoothly with constant sector ranks. Because a chart change acts by congruence, $\tilde{G} = J^\top G J$, absolute eigenvalue thresholds are chart dependent; a chart-invariant statement uses the generalized eigenvalue problem $G_i(c) v = \lambda g_0(c) v$ against a declared fiducial base tensor $g_0$ that transforms by the same congruence. We take $g_0$ to be the flat chart metric of the working framework (Section~\ref{sec:working_framework}). It is a declared computational background, not a physical metric, and the sectors are defined relative to it. A background-free alternative would use the consensus metric of Eq.~\eqref{eq:consensus_metric_covariant} in place of $g_0$ where that metric is positive definite; it degenerates wherever the agents' sections share a kernel, and we do not develop it.

An operational definition would replace the stipulated threshold by the observation channel, but it needs a declared map from $\mathcal{C}$ to the latent space. The natural choice through the means gives, for a Gaussian likelihood of precision $\Lambda_o$, the registered-information tensor $J^{(o)}(c) = d\mu_i^\top \Lambda_o d\mu_i$. This tensor is blind to covariance variation and degenerate wherever $d\mu_i$ has a kernel, so the pencil $G_i v = \lambda J^{(o)} v$ is generically singular, with infinite eigenvalues on directions in $\ker J^{(o)}$ outside $\ker G_i$. A well-posed version assigns $\ker J^{(o)}$ to the unregistered sectors and applies the pencil on a $g_0$-orthogonal complement of $\ker J^{(o)}$, so that observable directions are those along which belief variation is registered at the rate it occurs. No spatial likelihood channel is specified in this paper, so we keep the threshold form.

Register entry P6 postulates, for human cognitive agents, an observable sector with $|\mathcal{D}_{\text{obs}}| = 4$, one temporal and three spatial directions, with $|\mathcal{D}_{\text{obs}}| \ll |\mathcal{D}_{\text{subthresh}}| \ll |\mathcal{D}_{\text{internal}}|$. Since $\mathrm{rank} G_i \le \dim\mathcal{C}$, the postulate requires $\dim\mathcal{C} \ge 4$, and the stated hierarchy requires $\dim\mathcal{C}$ much larger. Neither holds on the working base $\mathcal{C} = \mathbb{R}^2$, where at most two sector eigenvalues are nonzero, so the postulate has no instance in this paper. The tensor $G_i$ is positive semidefinite, so nothing in its spectrum distinguishes a temporal direction; a Lorentzian reading needs register entry P4 or P5 (Section~\ref{sec:signature_resolution}). The framework neither establishes that the observable distribution is Frobenius-integrable, which a $(1+3)$ foliation would require, nor makes a quantitative prediction about measured spacetime.

Two senses of invisibility operate here and should not be conflated. At the base level, only the observable eigendirections of $G_i(c)$ are read as phenomenal spacetime, the remaining $\dim\mathcal{C} - |\mathcal{D}_{\text{obs}}|$ directions carrying subthreshold or internal eigenvalues. At the fiber level, the differential $d\sigma_i|_c: T_c\mathcal{C} \to T_{q_i(c)}\mathcal{B}$ samples at most $\dim\mathcal{C}$ directions of the $K_q(K_q+3)/2$-dimensional Gaussian belief manifold, so the unsampled fiber directions contribute nothing to $G_i(c)$ whatever the threshold.

\subsection{Observer-Dependent Reality}
\label{sec:observer}

Different agents pull back different metrics, $G_i \neq G_j$, over the same base. With the covariant Fisher block this difference has one source only, the difference between the agents' frame-stripped beliefs; with the ordinary fixed-gauge pullback, part of $G_i - G_j$ is gauge artifact, since a local frame change adds frame-gradient terms. Observer-dependence of the perceived geometry is therefore a statement about belief, not about gauge. It is not solipsism: agents remain coupled through the attention weights $\beta_{ij} \propto \exp[-D_{\mathrm{KL}}(q_i \| \Omega_{ij}[q_j])]$, and at a consensus stationary point their covariant pullbacks agree to the extent that their beliefs do. The resulting position combines within-frame validity with an intersubjective consensus that holds within a population of similar agents, not absolutely (status S).

The framework is classical and admits a global frame in Regime~I, so it models observer-dependent representation and belief, not the observer-dependent facts at issue in Brukner's no-go theorem~\cite{Brukner2018} and the Frauchiger-Renner argument~\cite{frauchiger2018quantum}. Its relation to Brukner's theorem, to the cross-perspective-links question of relational quantum mechanics~\cite{AdlamRovelli2022}, to QBism~\cite{Fuchs2017, DeBrotaFuchsSchack2020}, and to M\"uller's observer-first program~\cite{Mueller2020} is set out in Section~\ref{sec:discussion}.

``No view from nowhere'' has a limited formal meaning here. The base $\mathcal{C}$ is a presupposed smooth index manifold carrying no physical metric: it carries a background density and a fiducial chart metric, used for integration and for the sector thresholds above, and the metrics that represent perceived geometry are the induced pullbacks. Its dimension, topology, and smooth structure are inputs, so the construction derives metric data only and keeps the continuum that Wheeler's ``no continuum'' rejects~\cite{Wheeler1990}. Calling $\mathcal{C}$ noumenal is an analogy only (Section~\ref{sec:discussion}).

\subsection{Physical Quantities as a Derivation-Target Program}
\label{sec:phenomenological_interpretation}

Both induced tensors are dimensionless (Section~\ref{sec:pullback}). A physical line element requires an external length scale $L_0$, and a Lorentzian reading requires in addition a constant conversion speed, which the algebraic signature route does not supply and the causal route postulates without units (register entry P5).

Register entry P7 postulates that information, measured in nats by the free energy and the relative entropies, is the dimensionally fundamental quantity, and that physical dimensions are labels agents construct through consensus. The postulate reframes the dimensional question; it does not answer it, since every dimensionful constant, beginning with $L_0$, remains an external input that no construction in this paper produces. On the interpretive reading, information geometry on the fibers is the structure available to agents and physical measurements are the labels they assign. The falsifiability cost of the commitment is assessed in Section~\ref{sec:gauge_invariance_consensus}. The structural parallels to physical concepts (time from information updates, stiffness read as mass under the kinetic-metric postulate, action from the free-energy functional) are formal analogies with status S, summarized in Table~\ref{tab:physical_correspondences}.

\begin{table}[h]
\centering
\begin{tabularx}{\linewidth}{|l|X|X|}
\hline
\textbf{Physical Quantity} & \textbf{Framework Analog} & \textbf{Missing ingredient} \\
\hline
Time $t$ & Information clock $\vartheta_i$ & A presupposed ordering parameter is not removed (register entry P1); indefinite-signature extension (Section~\ref{sec:signature_resolution}) \\
Mass $m$ & Stiffness (Hessian of $\mathcal{F}$) under the kinetic-metric postulate & A second-order dynamics in which inertia and stiffness are independent; the flow used here is first order, with the Fisher metric as friction tensor (Section~\ref{sec:mass}) \\
Energy $E$ & Free energy density & An energy scale and an identification of $\mathcal{F}$ with a physical energy \\
Entropy $S$ & Information entropy $H[q]$ & Identification of $q$ with the physical ensemble, Boltzmann's constant $k_B$, and a phase-space measure~\cite{Jaynes1957}; the Gaussian $H[q]$ is a coordinate-dependent differential entropy \\
Action & Free energy functional & Legendre structure (stationarity correspondence only) \\
$\hbar$ & Minimal action per bit & Independently defined minimal action per bit \\
$c$ & Information propagation speed & Independently defined maximal epistemic speed with units (register entry P5) \\
\hline
\end{tabularx}
\caption{Correspondences between physical quantities and information-theoretic structures (status S), each labeled by the ingredient whose supply would upgrade it from analogy to derivation.}
\label{tab:physical_correspondences}
\end{table}

\section{Speculative Physical Readings: Gravity, Measurement, and Macroscopic Objects}\label{sec:speculative_physics}

\subsection{A Pullback-Curvature Reading of Gravity}
\label{sec:gravity_pullback}

By the three-tier reading of Section~\ref{sec:three_tiers}, the geometry an agent experiences as ambient is the structural pullback $G^{(s)}_i$ on the model fiber, and register entry P8 reads gravitational effects, if any, as curvature of this metric (status S). One exact statement holds: for constant isotropic model covariance $\Sigma_s = \sigma_s^2 I$ the structural metric is $G^{(s)} = \sigma_s^{-2} h$, with $h$ the first fundamental form of the immersion $\mu_s: \mathcal{C} \to \mathbb{R}^{K_m}$, and a constant conformal rescaling leaves the Levi-Civita connection unchanged and multiplies sectional curvature by $\sigma_s^2$, so the curvature is the intrinsic curvature of the immersion, given by the Gauss equation, rescaled by the model variance. Nothing further is derived: there is no field equation, source coupling, Newtonian limit, or gravitational constant, the structural metric is Riemannian, and no gravitational mass is defined, so the equivalence principle cannot yet be posed, and the reading is a placeholder, not a model of gravity. By contrast, Jacobson derives the Einstein equation as an equation of state from the Clausius relation on local Rindler horizons, with the entropy density fixing Newton's constant~\cite{Jacobson1995}, and Carlip sets out the obstructions that any emergent-gravity proposal must meet~\cite{Carlip2014}.

\subsection{Macroscopic Objects as Precision Attractors}

On this reading everyday objectivity has a classical mechanism. A high-precision agent, such as a macroscopic object with prior covariance $\Sigma_{\mathrm{ob}} \approx \delta I$ for small dimensionless $\delta$ (the precision scale of Section~\ref{sec:mass}), enters each coupling $\beta_{i,\mathrm{ob}} D_{\mathrm{KL}}(q_i \| \Omega_{i,\mathrm{ob}}[q_{\mathrm{ob}}])$ with such large precision that, where the coupling dominates, the alignment term pulls every observer's belief toward the object's transported belief, so high-precision systems anchor all observers. The anchoring is one-sided in a limited sense. Appendix~\ref{app:examples} shows that the back-action of the observers on the object's mean is suppressed by precision dominance, of order $\delta$ in absolute terms and $\delta^{1/2}$ in the object's Fisher metric, whereas the back-action on its precision is not: at stationarity the object's precision is reduced by a factor that tends to $1/(1+\beta)$ as $\delta \to 0$ and stays away from one at every $\delta$. The precision $\delta$ carries no physical units, and the dimensional bridge is the open question of Section~\ref{sec:phenomenological_interpretation}. Macroscopic objects are not modeled as agents anywhere in this paper, so the description extends the meta-agent construction by analogy (status S). Its structure, objectivity arising because many observers couple to the same high-precision record, is the classical counterpart of quantum Darwinism~\cite{Zurek2009}.

\subsection{An Inferential-Consensus Analogy for Measurement}
\label{sec:measurement_analogy}

Measurement can be read as the reverse of the macroscopic-object case: observers who hold different priors about a system couple to a high-precision apparatus, which dominates the free energy and pulls them into agreement. Classically this mechanism is precision-weighted opinion pooling~\cite{degroot1974, Aumann1976}, and its closest quantum-foundations relative is the quantum Darwinism noted above~\cite{Zurek2009}; environment-induced superselection~\cite{Zurek2003} needs a basis-selection structure that a framework without a Hilbert space lacks. The analogy is therefore suggestive only (status S): with no Hilbert space, Born rule, or superposition, the framework does not address the measurement problem. The closest inference-route constructions are Caticha's entropic dynamics, which derives Schr\"odinger evolution with the Born rule built into the representation~\cite{Caticha2015, Caticha2019, JohnsonCaticha2011}, and Reginatto's minimum-Fisher-information derivation, in which adding the Fisher term to the classical Hamilton-Jacobi action yields the Schr\"odinger equation~\cite{reginatto1998derivation}. We adopt Fisher information only as the metric that prices belief change, not Frieden's program of deriving physical law from extreme physical information~\cite{frieden1998physics}, whose claims of universality Lavis and Streater criticize~\cite{LavisStreater2002}. Whether anything resembling quantum measurement can arise awaits the density-operator fiber of item C5 (Section~\ref{sec:open_problems}).

\section{Consciousness: An Interpretive Stance}\label{sec:consciousness}

Every statement in this section has status S or C (Table~\ref{tab:claim_status}). No equation or proposition elsewhere depends on it, and a reader who rejects it keeps the mathematics intact. It records the author's interpretive stance so that its costs are visible.

\subsection{A Candidate for the Cognitive-Frame Transformation Law}
\label{sec:lahav_convergence}

Lahav and Neemeh write the cognitive-frame transformation law only for functionally identical cognitive systems and leave the law between distinct frames unspecified~\cite{LahavNeemeh2022, LahavNeemeh2025}. Section~\ref{sec:cognitive_reference_frames} proposes an explicit candidate, with the frame $U_i$ as cognitive frame of reference, the transport $\Omega_{ij}$ as the transformation between frames, and $D_{\mathrm{KL}}(q_i \| \Omega_{ij}[q_j])$ as the frame-invariant residual. The match is one of structural function, and its scope is narrow: at one base point Regime~I frames are removable (Proposition~\ref{prop:regime1_removability}), so a frame carries content only relative to fields held fixed, such as the observation channel. The correspondence departs from Lahav and Neemeh in two respects. They are relational physicalists, for whom physical reality is observer-relative but ontologically prior, whereas the stance here denies an observer-independent substrate beneath the agents; and neither account supplies a second postulate analogous to the constancy of the speed of light, since the signature constructions of Section~\ref{sec:worked_signature} are postulated, not selected. The qualia-indeterminacy worry that the relational reading dissolves is relocated, not removed: the only candidates for a frame-independent label are the invariants of Section~\ref{sec:consensus_metric}.

\subsection{Qualia, the Hard Problem, and Meta-Agents}

On the interpretive reading of register entry P9, the formal correlate of phenomenal character is the geometry of an agent's covariant pullback metric $G_i$. That tensor depends on the agent's belief section, not on its frame, so an inverted spectrum modeled as a passive re-expression of one agent's frame~\cite{Shoemaker1982} changes no invariant quantity; the formalism represents such an inversion as no difference at all. The stance relocates Chalmers' hard problem~\cite{Chalmers1995, Chalmers1996} without solving it: the explanatory gap moves from why neural computation feels like anything to why the geometry of inference does, a strategy shared with Russellian neutral monism~\cite{Russell1927, Strawson2006}, panpsychism~\cite{Goff2017, Chalmers2013consciousness}, and the identity claim of integrated information theory~\cite{Tononi2008, Tononi2015}, to which the framework adds vocabulary, not an explanatory mechanism.

The author's position in the panpsychism taxonomy, stated in the interpretive stance of Section~\ref{sec:cognitive_first}, is panprotopsychist in Chalmers' sense~\cite{Chalmers2013consciousness, Chalmers2016}: agents at level zero carry frame structure and take part in variational dynamics, and no phenomenal property is ascribed to them. Strawson's argument against brute emergence~\cite{Strawson2006} presses on constitution from protophenomenal structure as well, and the stance accepts the resulting combination problem as its central open question (status C).

The one conjecture offered is that conscious systems are meta-agents with active cross-level coupling that represent their constituents. It is a necessary condition that does not discriminate: meta-agent formation is triggered by coherence alone, so a highly ordered crystal qualifies, and the representation requirement is layered on the formation rule, not derived from it. The conjecture therefore does not compete with the integrated-information measure $\Phi$~\cite{Tononi2004, Tononi2016} as a test, and the framework draws no boundary between conscious and non-conscious systems.

\section{Gauge Invariance as Consensus and the Philosophy of Science}\label{sec:gauge_invariance_consensus}

The thesis of this section, that the gauge structure observed in physics reflects human cognitive architecture, has status C and may be unfalsifiable: any observed physics can be reinterpreted as reflecting cognitive constraints, and any disappearance of gauge structure as a change in cognitive architecture. The formalism is consistent with the thesis; it does not establish it.

\subsection{Model Alignment, Covariance, and the Selection Criterion}

The model-alignment term $\mathcal{F}_{\text{model}} = \sum_{ij} \int_{\mathcal{C}} \gamma_{ij}(x) D_{\mathrm{KL}}(s_i(x) \| \tilde{\Omega}_{ij}[s_j](x)) dx$ drives agents toward shared generative models, and consensus formation requires that scientific models be statable frame-independently, so that their content survives transport under $\tilde{\Omega}_{ij}$ (register entry P10); smallness of a pairwise divergence does not by itself establish a covariant transformation law. A century-old objection bounds what such a criterion can deliver. Kretschmann argued in 1917 that general covariance alone carries no empirical content, since any theory can be rewritten in generally covariant form by introducing auxiliary geometric objects~\cite{Kretschmann1917}, and the literature surveyed by Norton~\cite{Norton1993} has sharpened the distinction between physically substantive and merely notational covariance, and between coordinate invariance and gauge invariance proper, the latter constraining the form of the action, for example by forbidding explicit gauge-boson mass terms and fixing the universality of gauge couplings.
\begin{remark}[Kretschmann bound]\label{rem:kretschmann}
Cognitive shareability, as an intersubjectivity constraint on the form of theoretical statements, at most motivates general covariance and does not deliver gauge invariance proper; it is a candidate constraint on the formal scaffolding within which theories are stated, not an explanation of why nature exhibits the particular gauge groups it does.
\end{remark}
The $\mathrm{GL}(K_q)$ invariance of the divergence is already present for a single agent, so it is a consistency requirement that consensus must respect, not a symmetry that consensus generates. No Standard-Model factor is derived: $\mathrm{GL}(K_q, \mathbb{C})$ contains $\mathrm{U}(1) \times \mathrm{SU}(2) \times \mathrm{SU}(3)$, but every compact Lie group has a faithful finite-dimensional unitary representation and so embeds in some $\mathrm{U}(K_q)$, and the containment carries no selection content. The gauge group of physics is an empirical input, and whether free-energy minimization dynamically selects any subgroup is open. The historical progression from Galilean to Lorentz to general covariance to gauge invariance is consistent with the candidate cognitive-shareability constraint, which offers an interpretive reading of that history, not a derivation of physical law.

\subsection{Altered States, Counterfactual Architectures, and Informational Compatibility}

The cognitive reading re-describes altered states of consciousness as temporary shifts of one agent's frame relative to the consensus, so that the transport $\tilde{\Omega}_{\text{altered,normal}}$ lies far from the identity and the residual disagreement is large; this is a frame mismatch, not curvature, since $F = 0$ in Regime~I~\cite{Carhart-Harris2014, Swanson2018, CarhartHarrisFriston2019}. The consistency is not confirmatory (status S). Counterfactually, a population whose frames shifted coherently would pull back different metrics and construct different physics, the predictive-processing reading of perception~\cite{Clark2016, Seth2021, Hoffman2019, hohwy2013predictive} made geometric; this is a thought experiment, not a prediction, since no population with a substantially different cognitive architecture is accessible. Read this way, physics is a constitutive but revisable constraint on shareable description, in Friedman's relativized sense of the a priori~\cite{Friedman2001}; what the Discussion disavows is the apodictic a priori, fixed and universal, not this constitutive role. The resulting objectivity is that of robust consensus constrained by sensory evidence, and whether this reading explains or merely relabels remains open, since a realist account on which nature is simply uniform accommodates the same data at least as well.

\subsection{Scientific Knowledge as Collective Free-Energy Minimization}

On this reading scientific success is collective free-energy minimization under several constraints: fitting the empirical term $-\mathbb{E}_q[\log p(o \mid x)]$, reaching intersubjective consensus through the alignment term, and keeping the covariance needed for shareability. This is a pragmatic reconception in the spirit of constructive empiricism~\cite{vanFraassen1980} and of Laudan's problem-solving account of scientific rationality~\cite{Laudan1977}. It applies reflexively: the acceptance of this manuscript would itself be consensus-optimal free-energy minimization within a research community, and we offer the framework as a stance in van Fraassen's sense~\cite{vanFraassen2002}.

Kuhnian revolutions~\cite{Kuhn1962} read as collective frame changes $\phi_{\text{old}} \to \phi_{\text{new}}$ with accompanying model shifts $s_{\text{old}} \to s_{\text{new}}$. A uniform change of frame is a symmetry of $\mathcal{F}$, so the incommensurability analog lives in the model shift: the global transport $\Omega_{\text{old,new}}$ stays invertible, compatible with the local incommensurability of Hoyningen-Huene~\cite{Hoyningen-Huene1993}, and the formal content is communication breakdown, the collapse of inter-community attention when transported old-paradigm beliefs sit far from the new model's free-energy minima. Quinean underdetermination~\cite{Quine1951, Quine1975} is engaged in its weaker form: models of equal empirical adequacy, $D_{\mathrm{KL}}(p_{\mathrm{emp}} \| p_{s_1}) = D_{\mathrm{KL}}(p_{\mathrm{emp}} \| p_{s_2})$, are distinguished by further free-energy criteria (alignment, simplicity, unification), which break ties pragmatically. Whether such criteria are evidential~\cite{LaudanLeplin1991} or pragmatic~\cite{vanFraassen1980} is left open, and the framework is compatible with scientific pluralism~\cite{Longino2002}.

The framework implies epistemic limits of a structural kind. The base $\mathcal{C}$ is a presupposed smooth index manifold (Definition~\ref{def:base_manifold}) carrying no physical metric; every metric feature that represents perceived geometry is agent-induced through the pullback, so the formalism contains no agent-independent geometry to be known, and what is achievable is intersubjective agreement within an aligned community (Section~\ref{sec:observer}). The presupposed continuum is a departure from Wheeler, who answers regress with a loop, not a floor~\cite{Wheeler1990}. Wigner's unreasonable effectiveness~\cite{Wigner1960} is read as effective physics being the physics that achieves stable consensus, and Nagel's view from nowhere~\cite{Nagel1986} is the unattainable limit. What it would take to convert this structure into physics is the subject of Section~\ref{sec:open_problems}.

\section{Critical Open Problems and Future Directions}\label{sec:open_problems}

Several problems must be resolved before the framework can claim to explain physical reality, not merely supply structural analogies. Taken constructively, they form an ordered program: five mathematical problems, C1 through C5, and one derivation target, C6, each stated with what it would deliver and what it depends on, followed by open items that do not gate the program.

\begin{description}
\item[C1 (estimating loop holonomy).] At a single base point the independent links on the interaction graph of Section~\ref{sec:discrete_regime_ii} are the general case, and the frame-coboundary links of Regime~I are the graph-flat special case. A normalized population record whose law separates two link configurations that differ only in the conjugacy class of a loop holonomy exists~\cite[Ch.~10]{DennisGaugeVFE2026} (the operational loop-holonomy theorem), so detectability is settled at the level of the model. The mathematical problem is to construct a consistent estimator of loop-holonomy conjugacy classes from finitely many normalized population records under a stated noise model, with a finite-sample error bound. Its empirical application to a named target system would further require a held-out statistic and a tolerance fixed before the data, together with a control arm in which the links are frozen at a flat or random value. Under Regime~I every loop holonomy is trivial, so no such test has content there. A base-connection version, with holonomy along base loops, is separate and deferred~\cite{DennisGaugeVFESuppB2026}.
\item[C2 (consensus geometry).] Section~\ref{sec:consensus_metric} shows that the presence-weighted average of covariant Fisher pullbacks is gauge invariant without orbit averaging. What remains open is quantitative: a bound on the residual disagreement $G_i - G_j$ among similar agents in terms of their transported divergences at a consensus stationary point, which would give the intersubjective objectivity of Section~\ref{sec:observer} quantitative content. A candidate mechanism is the slow-channel alignment $\gamma_{ij} D_{\mathrm{KL}}(s_i \| \tilde\Omega_{ij} s_j)$, with each agent's prior inheriting alignment through the lagged predictive prior of Section~\ref{sec:cross_scale_shadows}. A regulated orbit average is needed only for the frame-twist block (Proposition~\ref{prop:regulator_obstruction}), and only if the signature constructions of C3 and C4 are kept.
\item[C3 (the four-generator construction beyond the commuting case).] Section~\ref{sec:signature_resolution} exhibits an all-real frame-twist form of signature $(-,+,+,+)$ from one skew and three commuting diagonal generators. The signature is fixed by that choice (register entry P4), so its availability is immediate from the indefinite trace form on $\mathfrak{gl}(K_q, \mathbb{R})$, and the construction is coordinate-flat; in Regime~I the frame-twist block is also pure gauge (Proposition~\ref{prop:regulator_obstruction}), so the construction is a property of a declared reference section. This item asks for a curved four-dimensional candidate beyond the abelian-exact regime, with quantitative control of the non-abelian corrections, built on a gauge-invariant form. The natural candidate, the relational trace form on agent pairs built from $\Omega_{ij}^{-1}\partial_\mu\Omega_{ij}$ (Section~\ref{sec:signature_resolution}), needs base derivatives of frames and is deferred~\cite{DennisGaugeVFESuppB2026}. The Lorentzian signature problem remains the central physical obstacle: Fisher-Rao geometry is Riemannian, and no tensor in this paper is shown to be both Lorentzian and the carrier of the clock and the sectors.
\item[C4 (a free-energy selection mechanism for the $1{+}3$ split).] Exhibit a mechanism by which free-energy dynamics compare frame configurations and select one compact and three non-compact commuting frame directions, the sharpened open problem of Section~\ref{sec:signature_resolution}. This item depends on C3 for the configuration space over which the selection would operate.
\item[C5 (the density-operator level-0 fiber).] Wheeler's first question, how come the quantum, is outside the scope of this paper, and this item names the first object a quantum extension would need. It would replace the real Gaussian belief fiber by a non-commutative statistical object, such as density matrices with a quantum relative entropy, for which the $\mathrm{GL}(K_q, \mathbb{C})$ enlargement of Section~\ref{sec:signature_resolution} supplies mathematical tools but not a quantum theory. It would further require inter-agent unitary transport with attention coupling under the subgroup $\mathrm{U}(K_q) \subset \mathrm{GL}(K_q, \mathbb{C})$ that preserves the Hermitian structure, with tensor-product bundles carrying entangled beliefs, a demonstration that the higher levels reduce to the classical Gaussian dynamics used here, and a choice of monotone (Petz-family) metric on the density-operator manifold~\cite{Petz1996}. The metric and transport tools exist in the quantum-information literature, and a free-energy principle for generic quantum systems has been proposed~\cite{Fields2022}; the open part is the multi-agent coupling.
\item[C6 (a dimensionless-ratio derivation).] Derive at least one dimensionless ratio of physical constants from information geometry alone, for example the fine-structure constant at a stated momentum-transfer scale (it runs from about $1/137$ in the Thomson limit to about $1/128$ at the $Z$ mass). Such a test is symmetric only under advance specification: the statistical manifolds, the comparison connection, and the \v{C}encov scale fixings, each unique only up to a positive scalar so that predictions must be scale-free ratios, must be fixed before the measured value is consulted, together with a completion criterion that certifies the construction independently of whether the derivation succeeds. No such criterion is available, so C6 is a derivation target, not a decidable problem, and failure under an incomplete or post-hoc construction is non-diagnostic. It depends on C2, because the hierarchy across which dimensionful quantities would be fixed presupposes a consensus geometry.
\end{description}

Three clusters of open items sit outside the ordered program. First, physical content. Matter fields and fundamental interactions are absent, and the construction makes no quantitative prediction for particle masses or coupling strengths; whether they could enter as excitations or topological features of the induced geometry is open, and the gravitational reading remains the postulate P8 (Section~\ref{sec:gravity_pullback}). Two natural extensions are recorded without development. One is a cross-level field strength built from the variation of frames across the level index $\ell$ as well as across base position; because coarse-graining steps are not invertible and compose as a one-parameter semigroup indexed by $\ell$, such an object would pair that semigroup with the structure group $G$, and no construction of it is attempted. The other is a pair of cross-bundle morphisms $\Phi: \mathcal{E}_{\mathrm{model}} \to \mathcal{E}_{\mathrm{state}}$ and $\tilde{\Phi}: \mathcal{E}_{\mathrm{state}} \to \mathcal{E}_{\mathrm{model}}$, read as the instantiation of a model in a state and as memory consolidation, which the companion theory defines at one base point~\cite[Ch.~2]{DennisGaugeVFE2026} but which do not enter the free-energy functional~\eqref{eq:free_energy_functional_final}.

Second, the scale question. No result on how the hierarchy behaves with growing agent count or depth is available; the exact coarse-graining statements and the nested bound (Table~\ref{tab:claim_status}) hold at fixed finite size. Its precise statements are the companion theory's open problems on a persistent multiscale hierarchy~\cite[Ch.~12]{DennisGaugeVFE2026} and on a family in which attention, learning, or holonomy selects the components of a collective~\cite[Ch.~13]{DennisGaugeVFE2026}. The one decidable test of the participatory reading of the formalism, the companion theory's pairing-selection conjecture~\cite[Ch.~14]{DennisGaugeVFE2026}, is stated in Section~\ref{sec:scope_limitations}. A finite-coupling transition remains conjectural and requires the prospective study of Section~\ref{sec:participatory}, including its frozen-attention-prior arm.

Third, open mathematical items pointed to from the body. On the noncompact $\mathrm{GL}^+(K_q)$ the frame mean of Eq.~\eqref{eq:meta_agent_frame_barycenter} needs a declared convention. Under any left-invariant Riemannian metric the Riemannian center of mass exists and is unique below Afsari's radius but is only left-equivariant~\cite{Afsari2011}, while the exponential barycenter of the canonical Cartan-Schouten connection is left-, right-, and inversion-equivariant and exists uniquely for sufficiently concentrated frames~\cite{PennecArsigny2013}; the additive form $\exp(\sum_i w_i \phi_i)$ is the first fixed-point iterate of the latter from the identity. What is open is the choice of convention, which Proposition~\ref{prop:parent_frame_gauge} (Section~\ref{sec:meta_agent_definition}) identifies as a gauge choice. The kinetic-metric postulate of Section~\ref{sec:velocity_quadratic} needs its own second-order equation of motion before any frequency test can be posed, since the dynamics used in this paper are first order, with the Fisher metric as friction tensor. A separate timescale for each level of the hierarchy (Section~\ref{sec:ouroboros}) is not constructed.

\section{Conclusion}\label{sec:global_conclusion}

This paper has given a classical, multi-agent, information-geometric model of parts of Wheeler's participatory program, and Table~\ref{tab:wheeler_correspondence} records the status of each part. For registered change as a clock (W1), each agent's clock is the Fisher length of its belief history, with an exact dissipation identity along an autonomous single-rate natural-gradient flow, and a binary registration channel gives Wheeler's bit its meaning as at most one bit of mutual information per yes-no registration (Section~\ref{sec:time_info}, Proposition~\ref{prop:binary_registration}). For observer-dependent geometry (W2), metrics on a base that carries no physical metric are pulled back from belief sections; with the covariant derivative the pullback is gauge invariant, so its observer-dependence is dependence on belief (Sections~\ref{sec:it_from_bit} and~\ref{sec:consensus_metric}). For one world out of many observer-participants (W3), intersubjective agreement is a stationary point of an explicit consensus functional, a formal counterpart of Wheeler's adoption of F{\o}llesdal's dictum, with the residual transported divergence as the frame-independent disagreement (Section~\ref{sec:collective_geometry}). For participation across scales (W4), exact coarse-graining builds parents from their members, and a parent's belief returns to its members as a one-episode-lagged predictive prior; the hierarchy closes at an energy-level apex that is, in any generative reading, a declared root, the formal residue of Wheeler's refusal of a tower of turtles (Sections~\ref{sec:cross_scale_shadows} and~\ref{sec:ouroboros}). These are constructions and derivations under stated assumptions; their identification with Wheeler's notions has status S.

What the model does not deliver (W5) is as definite. It contains no discrete irreversible registration with complementarity, so it does not model Wheeler's elementary quantum phenomenon; its functional, belief family, and divergence are fixed in advance, so it does not model law without law; it presupposes a smooth index manifold and an ordering parameter, so it keeps a continuum; its feedback relaxes to stationarity unless observations drive it, so it is not a self-excited circuit; and how come the quantum is outside its scope, with the density-operator fiber of C5 as the first object an extension would need. The signature constructions exhibit $(-,+,+,+)$ forms for postulated generators without selecting them, and no construction supplies physical units.

The companion theory~\cite{DennisGaugeVFE2026} proves the single-point mathematics in greater generality and declines, without refuting, each physical identification postulated in Part~II. Items C1 through C5 of Section~\ref{sec:open_problems} are mathematical problems and C6 is a derivation target. Solving them would decide questions about the formalism, not the participatory reading, whose one stated decidable test is the pairing-selection conjecture of the companion theory~\cite[Ch.~14]{DennisGaugeVFE2026}.
